\section{Conclusion}

This paper presented a modular closed-loop simulation framework
for vision-based autonomous spacecraft rendezvous and evaluated
each component of the perception, estimation, and guidance
pipeline through six targeted experiments.

The dynamics experiment confirmed that the HCW state-transition
implementation is numerically consistent with high-accuracy DOP853
integration to sub-nanometer position precision over three orbital
periods, validating the dynamics model as a suitable foundation
for the subsequent control and estimation studies.

The pose-estimation experiments established two complementary
properties of the visual pipeline.
First, Gauss--Newton nonlinear refinement substantially reduces
the raw EPnP translation error across all tested noise levels,
achieving a $76\%$ reduction at $5$~px pixel noise.
Second, RANSAC-based correspondence rejection is essential for
robustness: bare EPnP failed in up to $100\%$ of trials at
$30\%$ outlier contamination, while RANSAC--EPnP maintained
zero solver failures and sub-$8$~cm translation accuracy
throughout the tested range.

The state-estimation comparison showed that MEKF and UKF achieve
effectively identical translational performance under the
controlled Gaussian measurement model, with position and
velocity RMSE of $0.60$~m and $0.023$~m/s respectively.
The MEKF produced a marginally lower attitude RMSE
($1.33^\circ$ vs $1.34^\circ$), while the UKF required
approximately half the wall-clock execution time.

The guidance and control experiment revealed a clear performance
trade-off between the two strategies.
MPC achieved $58\%$ lower cumulative $\Delta v$ than LQR
($7.02$~m/s vs $16.83$~m/s) for the same rendezvous task,
but at a $300\times$ higher per-step computational cost
($3.00$~ms vs $0.01$~ms).
Both controllers converged to sub-millimetre final range,
confirming that terminal accuracy is not a differentiating factor
for the unconstrained scenario evaluated here.

The end-to-end Monte Carlo study demonstrated $100\%$ docking
success over 25 trials for all three configurations tested:
perfect-state feedback, direct noisy-measurement EKF, and the
full vision--MEKF pipeline.
The full pipeline achieved cumulative $\Delta v$ comparable to
the perfect-state reference ($\approx40.2$~m/s), substantially
below the direct EKF configuration ($\approx68.2$~m/s), showing
that RANSAC--EPnP visual measurements provide cleaner state
estimates than direct noisy position measurements under the
tested conditions.
The pixel-noise ablation further confirmed that $100\%$ docking
success is maintained from $0.3$~px to $4.0$~px, demonstrating
substantial tolerance of the closed-loop system to image-plane
measurement degradation.

Taken together, the results validate the proposed architecture
at both the subsystem and system level.
The principal limitation of the present study is its reliance on
a simplified synthetic camera model with parametric keypoints
and additive Gaussian noise; real imagery introduces
illumination variation, partial target visibility,
texture-dependent matching ambiguities, and hardware-specific
noise characteristics not captured by this model.

Future work should address these limitations through integration
of a higher-fidelity rendering pipeline, extension to
unstructured natural-feature keypoints, and hardware-in-the-loop
testing with a real camera and representative embedded processor.
Incorporating $J_2$ and atmospheric drag perturbations into
the dynamics model and evaluating the MPC approach-cone constraint
in the closed-loop end-to-end pipeline are also important
directions for future investigation.
